% Part: sets-functions-relations
% Chapter: relations
% Section: reflections
%
\documentclass[../../../include/open-logic-section]{subfiles}

\begin{document}

\olfileid{sfr}{rel}{ref}
\olsection{Filosofische beschouwingen}

In \olref[set]{sec} hebben we relaties gedefinieerd als bepaalde
verzamelingen. Het is goed om even stil te staan bij een filosofische
vraag: wat \emph{doet} zo'n definitie? Het is zeer de vraag of we
willen zeggen dat we daarmee hebben \emph{ontdekt} welke dingen
metafysisch gezien identiek zijn; bijvoorbeeld dat de ordeningsrelatie op $\Nat$
\emph{bleek} samen te vallen met de verzameling $R=
\Setabs{\tuple{n,m}}{n, m \in \Nat\text{ en } n < m}$ die we in
\olref[set]{sec} hebben gedefinieerd. Daarvoor zijn drie redenen.

Ten eerste hebben we in \olref[set][pai]{wienerkuratowski} gedefinieerd
dat $\tuple{a, b} = \{\{a\}, \{a, b\}\}$. Beschouw in plaats daarvan de
definitie $\lVert a, b\rVert = \{\{b\}, \{a, b\}\} = \tuple{b,a}$. Als
$a \neq b$, dan geldt $\tuple{a, b} \neq \lVert a,b\rVert$. We hadden
echter net zo goed $\lVert a,b\rVert$ in plaats van $\tuple{a,b}$ als
definitie van een geordend paar kunnen kiezen. Beide definities zouden
even goed hebben gewerkt. Er zijn dus twee even geschikte
verzamelingen die als de ordeningsrelatie op de natuurlijke getallen
zouden kunnen gelden:
\begin{align*}
		R &= \Setabs{\tuple{n,m}}{n, m \in \Nat \text{ en }n < m}\\
		S &= \Setabs{\lVert n,m\rVert}{n, m \in \Nat \text{ en }n < m}.
\end{align*}
Omdat $R \neq S$, kunnen ze volgens het extensionaliteitsbeginsel niet
\emph{allebei} identiek zijn aan de ordeningsrelatie op~$\Nat$. Het zou
echter volstrekt willekeurig, en daardoor nogal ongemakkelijk, zijn om
te beweren dat $R$ in plaats van $S$ (of omgekeerd) feitelijk de
ordeningsrelatie \emph{is}. (Dit is een zeer eenvoudig geval van een
argument tegen het reduceren van zulke objecten tot verzamelingen, dat
Benacerraf in \citeyear{Benacerraf1965} bekend heeft gemaakt. We komen
er nog enkele keren op terug.)

Ten tweede: als we menen dat \emph{elke} relatie met een verzameling
moet worden geïdentificeerd, dan moet ook de elementrelatie zelf,
$\in$, een bepaalde verzameling zijn. Dat zou de verzameling
$\Setabs{\tuple{x,y}}{x \in y}$ moeten zijn. Maar bestaat die
verzameling? In het licht van Russells paradox is het allerminst
vanzelfsprekend dat ze bestaat. Sterker nog,
\oliflabeldef{cumul:::part}{de verzamelingenleer die we in
\olref[cumul][][]{part} ontwikkelen, zal het bestaan van deze
verzameling \emph{ontkennen}.\footnote{We lopen op de tekst vooruit;
dit is de reden. Neem ter tegenspraak aan dat $I =
\Setabs{\tuple{x,y}}{x \in y}$ bestaat. Dan is $\bigcup \bigcup I$ de
universele verzameling, in strijd met
\olref[sfr][z][sep]{thm:NoUniversalSet}.}}{het mogelijk is de
verzamelingenleer nauwkeurig als een axiomatische theorie op te bouwen,
en ook die theorie zal het bestaan van deze verzameling ontkennen.}
Ook als we sommige relaties dus als verzamelingen kunnen behandelen,
blijft de elementrelatie een bijzonder geval.

Ten derde: toen we relaties met verzamelingen ``identificeerden'',
spraken we af dat we $Rxy$ mochten schrijven voor $\tuple{x,y} \in R$.
Dat kan zolang we de elementrelatie, ``$\in$'', \emph{als} predicaat
behandelen. Maar als ``$\in$'' zelf een of andere verzameling aanduidt,
bestaat de uitdrukking ``$\tuple{x,y} \in R$'' slechts uit drie
singuliere termen die verzamelingen aanduiden: ``$\tuple{x,y}$'',
``$\in$'' en ``$R$''. Met zo'n rij namen kan men evenmin een uitspraak
uitdrukken als met de onzinreeks ``het kopje pennenhouder de tafel''.
Ook hier geldt dat sommige relaties weliswaar als verzamelingen kunnen
worden behandeld, maar dat de elementrelatie een bijzonder geval moet
blijven. (Hier worden een eenvoudige versie van Freges paradox van het
begrip \emph{paard} en een bekend bezwaar dat Wittgenstein ooit tegen
Russell inbracht, samengenomen.)

Wat volgt hieruit? Er is niets \emph{mis} mee om te zeggen dat relaties
op de getallen verzamelingen zijn. We moeten alleen begrijpen hoe die
uitspraak bedoeld is. We doen geen uitspraak over welke dingen
metafysisch gezien identiek zijn. We merken slechts op dat we in bepaalde contexten bepaalde
relaties als bepaalde verzamelingen kunnen en zullen \emph{behandelen}.
\end{document}